\subsection{A shared Gram lift before the fibre supremum (23 October)}
\label{sec:shared-response-gram-lift}

\paragraph{Scope.}
The following extraction is dimension-free and applies to every actual
disc-valued response pair. The stability constraint additionally requires a
globally near-maximizing pair, as in
\S\ref{sec:fresh-second-chaos-stability}. Neither the extraction nor the
floating-point diagnostics below gives a new numerical endpoint. In particular,
the atomic laws of \S\ref{sec:binding-audit-20261023} do not specify the new
observables and cannot be declared excluded by assigning them convenient values.

\paragraph{The common projection budget.}
Use the realification of complex Gaussian $L^2$, with scalar product
$\operatorname{Re}\mathbb E\overline u v$. Let $P$ be the full first-chaos
projection, and let $P_0$ project onto any chosen complex subspace of that
chaos. Fix bounded real features $A_1,\ldots,A_r,B_1,\ldots,B_s$ and form the
\emph{single} tuple
\[
 v=(fA_1,\ldots,fA_r,gB_1,\ldots,gB_s),\qquad
 x_i=P_0v_i,\quad z_i=v_i-x_i.
\]
Write $\operatorname{Gr}(u)_{ij}=\operatorname{Re}\langle u_i,u_j\rangle$.
The measured matrix $X=\operatorname{Gr}(x)$ and the residual matrices
\[
 V=\operatorname{Gr}(z)=\operatorname{Gr}(v)-X,\qquad
 K=\operatorname{Gr}(Pv-x)
\]
satisfy
\begin{equation}\label{eq:shared-gram-interval}
 \boxed{0\preceq K\preceq V,\qquad
       \operatorname{Gr}(Pv)=X+K.}
\end{equation}
Indeed $Pz_i=Pv_i-x_i$. For any real $c$,
$c^{\mathsf T}Kc=\|P\sum c_i z_i\|^2$ lies between zero and
$\|\sum c_i z_i\|^2=c^{\mathsf T}Vc$. Orthogonality to the measured
space proves the other two identities. No independence, conditional
Gaussianity, tensor-rank bound or estimate of an individual Hermite tail is used.

Crucially, the mixed block of $\operatorname{Gr}(v)$ is
\[
 \operatorname{Gr}(v)_{i,r+j}
   =\mathbb E\bigl[\operatorname{Re}(\overline f g)A_iB_j\bigr].
\]
It is a moment of the \emph{same pair}. Separate Bessel budgets discard it.
The dimension of \eqref{eq:shared-gram-interval} is $r+s$, regardless of the
number of unmeasured Gaussian directions.

\paragraph{A sharp scalar completion screen.}
For one feature on each branch, put $c=X_{12}$,
$v=V_{12}$ and $R=\sqrt{V_{11}V_{22}}$. Every possible residual coefficient
$k=K_{12}$ obeys the sharper interval
\begin{equation}\label{eq:shared-cross-interval}
 \boxed{|2k-v|\le R.}
\end{equation}
One proof uses the reflection $2P-I$: it is an isometry and
$2\langle Pz_1,Pz_2\rangle-\langle z_1,z_2\rangle
=\langle(2P-I)z_1,z_2\rangle$. Apply Cauchy--Schwarz.
For $V=I_2$, projection onto $(1,1)$ attains $k=1/2$; projection onto
$(1,-1)$ attains $k=-1/2$. Thus the radius $R/2$ cannot be reduced from
Gram information alone. The separate estimate $|k|\le R$ has twice that radius
when $v=0$.

Suppose stability gives $|c+k|\le a$. Eliminating \emph{all} completions yields
\begin{equation}\label{eq:shared-completion-screen}
 \boxed{|2c+v|\le R+2a.}
\end{equation}
This follows by subtracting $2k-v$ from $2(c+k)$ and using the triangle
inequality. For example, $c=1/5$, $v=0$, $R=1/10$ and $a=13/100$ fail
the screen: $2/5>9/25$. Separate Cauchy and stability still allow $k=-2/25$.
This is an exact example separating the two completion tests, not an example
at a Grothendieck binding node.

\paragraph{Connection to actual response stability.}
Retain the weights and measured moments before taking any fibre supremum:
\[
 D_f[i,j]=\mathbb E(|Tg|A_iA_j),\qquad
 D_g[i,j]=\mathbb E(|Tf|B_iB_j),\qquad
 C=(X+K)_{fg}.
\]
If the pointwise feature budgets are at most $L^2$, a globally
$\delta$-near-maximizing pair has the necessary finite-dimensional completion
\begin{equation}\label{eq:shared-stability-completion}
 0\preceq K\preceq V,\qquad
 \begin{pmatrix}D_f&\varepsilon(X+K)_{fg}\\
 \varepsilon(X+K)_{gf}&D_g\end{pmatrix}
       +10L^2\sqrt\delta\,I\succeq0.
\end{equation}
For scalar features, one may use
$a=\sqrt{(D_f+10L^2\sqrt\delta)(D_g+10L^2\sqrt\delta)}/\varepsilon$
in \eqref{eq:shared-completion-screen}. The signed weights
$\operatorname{Re}(\overline f Tg)$ and
$\operatorname{Re}(\overline g Tf)$ give the stronger original Hessian
constraint. They must not silently be substituted for the positive weights
in a computation that assumes the latter.

\paragraph{Finite dual prices for the residual block.}
For real symmetric matrices $Z,Y$, the two finite checks
$Y\succeq0$ and $Y-Z\succeq0$ imply
\begin{equation}\label{eq:shared-gram-price}
 \langle Z,K\rangle\le\langle Y,V\rangle
 \qquad(0\preceq K\preceq V).
\end{equation}
Here $\langle A,B\rangle=\operatorname{tr}(AB)$. Indeed
$\langle Z,K\rangle\le\langle Y,K\rangle\le\langle Y,V\rangle$.
Rational Gram factorizations of $Y$ and $Y-Z$ suffice to check the prices;
an optimizer's eigenvalues are unnecessary. The rank-one version is included
in the Lean module below.

This budget is exact for arbitrary Hilbert-space feature vectors: given
$0\preceq K\preceq V$, choose $A=V^{\dagger/2}KV^{\dagger/2}$ on
$\operatorname{ran}V$. Then $0\preceq A\preceq I$ and the block operator
\[
 \Pi=\begin{pmatrix}A&\sqrt{A(I-A)}\\
                      \sqrt{A(I-A)}&I-A\end{pmatrix}
\]
is an orthogonal projection. Applying it to the columns of
$(V^{1/2},0)^{\mathsf T}$ realizes $K$. This does \emph{not} realize
arbitrary such columns as $fA_i,gB_j$ on Gaussian space. The remaining
structure must enter through their local response moments.

\paragraph{A moment-constrained fibre support interface.}
There is a dimension-free way to choose the features so that their response
moments can be replaced by moments of the two polynomial images, with a
controlled near-maximizer error. This avoids declaring arbitrary bounded
responses independent of those images.

\paragraph{An image-determined near-maximizer lift.}
Let $a=|Tg|$, $b=|Tf|$, $\phi_g=\operatorname{phase}(Tg)$ and
$\phi_f=\operatorname{phase}(Tf)$. A globally $\delta$-near-maximizing pair
satisfies
\begin{equation}\label{eq:shared-weighted-alignment}
 \mathbb E\,a|f-\phi_g|^2\le2\delta,\qquad
 \mathbb E\,b|g-\phi_f|^2\le2\delta.
\end{equation}
For example, the first bound follows from the pointwise disc inequality
$a|f-\phi_g|^2\le2(a-\operatorname{Re}\overline f Tg)$ and
$\mathbb E|Tg|\le M$, since replacing $f$ by $\phi_g$ is admissible in
the same dimension. Both inequalities hold for the original pair; successive
phase replacements or existence of a maximizer are unnecessary.

Choose real bounded image-dependent rules $\alpha_i,\beta_j$ with
$\sum\alpha_i^2,\sum\beta_j^2\le1$, fix $\tau>0$, and set
\[
 A_i=\sqrt{\frac a{a+\tau}}\,\alpha_i,
 \qquad B_j=\sqrt{\frac b{b+\tau}}\,\beta_j.
\]
Then, summing over features,
\begin{equation}\label{eq:shared-cutoff-alignment}
 \sum_i\|(f-\phi_g)A_i\|_2^2\le\frac{2\delta}{\tau},\qquad
 \sum_j\|(g-\phi_f)B_j\|_2^2\le\frac{2\delta}{\tau}.
\end{equation}
Indeed $\sum A_i^2/a\le1/\tau$, including $a=0$ by continuity.
Define the image tuple
$\widetilde v=(\phi_g A,\phi_f B)$. Its full, measured and residual
Gram matrices differ in operator norm from those of $v=(fA,gB)$ by at most
\[
 4\sqrt{2\delta/\tau}.
\]
To see this, the synthesis operators of both tuples have Hilbert--Schmidt
norm at most $\sqrt2$, and their difference has norm at most
$2\sqrt{\delta/\tau}$. Use
$\|U^*U-\widetilde U^*\widetilde U\|
\le(\|U\|+\|\widetilde U\|)\|U-\widetilde U\|$.
The same argument applies after any orthogonal projection.
For the cross block alone the sharper error is
$2\sqrt{2\delta/\tau}$, because each individual branch has synthesis norm
at most one. Consequently the \emph{image-determined} necessary condition is
\begin{equation}\label{eq:shared-image-stability}
 \boxed{\begin{pmatrix}D_f&\varepsilon\widetilde C\\
       \varepsilon\widetilde C^{\mathsf T}&D_g\end{pmatrix}
       \succeq-\left(10+2\varepsilon\sqrt{2/\tau}\right)\sqrt\delta\,I,}
\end{equation}
where $\widetilde C=\operatorname{Gr}(P\widetilde v)_{fg}$ and $D_f,D_g$
are unchanged. Its measured and unmeasured matrices satisfy the exact common
Gram interval for $\widetilde v$. For fixed positive $\tau$, the error vanishes
uniformly in the dimension. This is a necessary condition on global
near-maximizers, not a bound on all arbitrary cubic perturbations.

At fixed polynomial-image state the features and $\widetilde v$ are now
known functions, rather than moments of free response variables. One possible
rule uses clipped $\operatorname{Re}(\overline{\phi_g}\ell_i)/(a+\eta)$,
normalized to the declared feature budget, and its other-branch counterpart.
The image law must still come from the same pair under $T$; arbitrary pairs
of polynomial images are an outer relaxation. No claim is made that the
present radial tables already enforce that compatibility.

\paragraph{The local moment support.}
Here is a way to retain those moments rather than append free matrices after
the old table. Fix a remote state $w$, including the conditional polynomial
images $(Tf)(\cdot,w),(Tg)(\cdot,w)$, and any bounded feature rule. For
fixed real quadratic-dual variable $y$, let $F_y(z;u,v)$ be the linear
two-response payoff from the old fibre identity and let
$m(z;u,v)$ collect the measured coefficients, the joint Gram entries and the
weighted feature moments just defined. If their fibre expectation is $b$, then
for \emph{every} real multiplier vector $\theta$,
\begin{align}
 \sup_{\substack{|f|,|g|\le1\\\mathbb E_z m(z;f,g)=b}}
 \{\mathbb E_zF_y(z;f,g)-y^{\mathsf T}My\}
 &\le\mathbb E_z\sigma_z(y,\theta)-y^{\mathsf T}My-\theta\cdot b,
 \label{eq:shared-local-support}\\
 \sigma_z(y,\theta)
 &=\sup_{|u|,|v|\le1}
       \{F_y(z;u,v)+\theta\cdot m(z;u,v)\}.\nonumber
\end{align}
This is just the pointwise support inequality, integrated before subtracting
the measured moments. Supremizing in $y$ preserves the bound; infimizing
over chosen prices remains valid. No tensor supremum is interchanged with an
infimum. Each inner support is on two complex discs, hence four real variables.
It retains $\operatorname{Re}(\overline u v)$; replacing it by two independent
one-disc supports would discard precisely the joint information.

For response-independent features the moment integrands give the quadratic
support
\[
 c+\sup_{|u|,|v|\le1}
 \{a|u|^2+b|v|^2+\operatorname{Re}(\overline p u+\overline q v)
                         +2\operatorname{Re}(\overline u wv)\}.
\]
One disc can be eliminated exactly. For $r\ge0$ define
\[
 \chi_b(r)=\begin{cases}b+r,&2b+r\ge0,\\
                         -r^2/(4b),&2b+r<0.\end{cases}
\]
Maximizing first the phase and then the modulus of $v$ gives the equivalent
\emph{two-real-variable} support
\begin{equation}\label{eq:shared-two-real-support}
 c+\sup_{|u|\le1}\{a|u|^2+\operatorname{Re}(\overline p u)
                   +\chi_b(|q+2\overline w u|)\}.
\end{equation}
Indeed $\chi_b(r)=\max_{0\le t\le1}(bt^2+rt)$; its maximizer is $1$ in
the first branch and $-r/(2b)$ in the second. Both the bound and attainment
are kernel-checked in the Lean module. Negative diagonal prices must use the
second branch when applicable: assuming optimal responses are unimodular
would otherwise give a false support value. If $a,b\ge0$, separate convexity
puts both responses on their circles and reduces the last search to one angle.
These simplifications do not apply unchanged to features depending nonlinearly
on the response, such as clipped response-weighted scores.

The standalone checker \code{tools/joint\_disk\_support.py} implements a
rational enclosure of \eqref{eq:shared-two-real-support}. On a rectangle $R$
meeting the unit disc, let $n_-$ be its minimum squared distance to the origin
and $n_+=\min\{1,\max_{\rm corners}|u|^2\}$. Bound the quadratic term by
$an_+$ if $a\ge0$, and by $an_-$ otherwise. The linear term is bounded by
the minimum of $|p|$ and its maximum over the four corners. Finally use
\[
 |q+2\overline w u|\le
 \min\{|q|+2|w|,\max_{\rm corners}|q+2\overline w u|\}.
\]
The corner bound follows from convexity of the norm; $\chi_b$ is increasing
because it is the supremum of $bt^2+rt$ with $t\ge0$. Adding these three
bounds gives a valid cell upper bound. Feasible rational points give lower
bounds. Rational square-root enclosures follow from integer square roots;
dyadic subdivision preserves a cover of the disc, and cells are discarded
only when disjoint from it or below the incumbent lower bound. Thus the
returned interval is rigorous even if a cell budget is exhausted before
reaching the requested tolerance. This checker covers fixed rational local
coefficients; it does not integrate over Gaussian states or certify a norm cap.

\paragraph{An exact $3\times3$ complex SDP for the local support.}
There is a stronger simplification of the same quadratic support. Put
\[
 H=\begin{pmatrix}
 0&\overline p/2&\overline q/2\\
 p/2&a&w\\q/2&\overline w&b
 \end{pmatrix}.
\]
Then, exactly,
\begin{align}
 \sup_{|u|,|v|\le1}(1,u,v)^*H(1,u,v)
 &=\max_{X\succeq0,\ X_{00}=1,\ X_{11},X_{22}\le1}
           \operatorname{tr}(HX)\label{eq:shared-exact-local-sdp}\\
 &=\min_{\substack{\lambda_1,\lambda_2\ge0\\
             \operatorname{diag}(\lambda_0,\lambda_1,\lambda_2)-H\succeq0}}
            (\lambda_0+\lambda_1+\lambda_2).\nonumber
\end{align}
Here $X$ is complex Hermitian and $\lambda_0$ is real without a sign
restriction. This is a local two-response result, not an exact SDP for the
whole Grothendieck problem.

Equivalently, the complete local quadratic moment body is
\[
 \boxed{\operatorname{conv}\{zz^*:z=(1,u,v),\ |u|,|v|\le1\}
       =\{X\succeq0:X_{00}=1,\ X_{11},X_{22}\le1\}.}
\]
Thus a loss in this part of a lifted certificate can be removed without a
higher moment hierarchy. This is local convexification only; it does not
identify which fields of such moments arise from a common Gaussian response
pair with prescribed polynomial projections.

For completeness, the feasible set of $X$ is compact and convex. If an extreme
point has rank $r\ge2$, write $X=RR^*$ with $R$ of full column rank. The real
space of Hermitian $r\times r$ matrices has dimension $r^2\ge4$. Therefore
there is a nonzero Hermitian $U$ with
$\operatorname{diag}(RUR^*)=0$, since this imposes at most three real equations.
For sufficiently small $t>0$, the distinct matrices
$R(I\pm tU)R^*$ are positive semidefinite and preserve all three diagonals.
Their average is $X$, contradicting extremality. Every extreme point thus has
rank one. Since $X_{00}=1$, write it as $(1,u,v)(1,u,v)^*$ after a common phase
rotation, with $|u|,|v|\le1$. A linear objective attains its optimum at an
extreme point, proving the first equality. The strictly feasible point
$\operatorname{diag}(1,1/2,1/2)$ and a sufficiently large diagonal dual matrix
give strong duality and attainment for the displayed primal and dual.

This is the usual complex correlation-matrix rank argument, here adapted to
the two-disc support; see Li--Tam, \emph{A note on extreme correlation matrices},
SIAM J. Matrix Anal. Appl. 15 (1994), 903--908,
\href{https://doi.org/10.1137/S0895479892240683}{doi:10.1137/S0895479892240683}.
The real symmetric perturbation space has dimension $r(r+1)/2$ instead;
the argument permits rank two with three real diagonal constraints. Thus this
exact local simplification specifically uses complex response phases. Adding
more response branches also invalidates the rank-one dimension count.

For rational data, any rational feasible dual prices are accepted by checking
the seven principal minors of the $3\times3$ Hermitian slack, together with
$\lambda_1,\lambda_2\ge0$. A floating SDP only proposes the prices; a small
rational diagonal loading makes the slack certifiably positive when necessary.
\code{tools/two\_disc\_sdp.py} contains a standard-library exact verifier and
a separate optional proposal mode. Its four local examples, including an
interior coupled optimizer, have exact dual--primal gaps below $10^{-8}$.
This replaces the rectangle mesh for the quadratic local support. The rectangle
checker remains an independent enclosure method. Neither checker supplies the
Gaussian integral or the global joint-image law cover.

There is also a direct continuous-cell rule. If
$H(s)=\sum_j\omega_j(s)H_j$, with nonnegative weights summing to one, interpolate
feasible dual prices by the same weights. The slack is the corresponding
convex combination of the certified positive semidefinite slacks, so the local
support is at most $\sum_j\omega_j(s)\sum_i\lambda_{j,i}$ throughout the cell.
This applies directly when the coefficient matrix is affine in the cell
parameters; more general inputs need a certified coefficient enclosure.
The old negative quadratic in $y$ remains separate and is spent once.

For a fixed lifted state the dual support is convex in the moment prices:
it is a supremum of affine functions of those prices, even after supremizing
in $y$. This does not make the inner maximization in $y$ concave or justify
exchanging the tensor supremum and price infimum.

For a positive definite $M$, the outer quadratic dual is coercive. Bounded
features and fixed polynomial images give an integrable bound
$C_0+C_1\|y\|-\lambda_{\min}(M)\|y\|^2$ on its value, with an added
explicit constant for the moment prices. Thus a finite ball in $y$ and the two
discs suffice for interval subdivision; Gaussian polynomial tails are paid as
in the old potential checker. This gives a finite support obligation for each
specified lifted state and price. It does \emph{not} yet give a finite cover
of all admissible remote laws of those states. In particular the conditional
images must remain the actual images of the same responses, rather than free
polynomials chosen separately for each moment.

\paragraph{A bilinear alternative for fixed feature rules.}
There is also a useful way to avoid the absolute-weight price when the real
features $A,B$ are response-independent. Let $M_A$ denote multiplication by
$A$ and use the signed Hessian form from
\S\ref{sec:fresh-second-chaos-stability}. For $\lambda\ge0$,
\[
 J(f,g)+\lambda H(A,B)
 =\operatorname{Re}\langle f,S_\lambda g\rangle,\qquad
 S_\lambda=T+\lambda(M_{A^2}T+TM_{B^2}+2\varepsilon M_APM_B).
\]
If the features have a uniform bound in every dimension, the near-maximizer
Hessian theorem and a maximizing sequence imply
\[
 \boxed{\|T\|_{\infty\to1}\le
        \sup_n\|S_\lambda\|_{\infty\to1}.}
\]
Indeed, at a $\delta$-near-maximizer the right payoff is at least
$\|T\|-\delta-10\lambda L^2\sqrt\delta$; now let $\delta\downarrow0$.
The same argument permits finite nonnegative sums of bounded feature tests.
This is a new certificate interface, not a claim that the modified norm is
smaller or already bounded. Response-dependent clipped scores do not define
one fixed bilinear operator and must instead use the lifted support interface.

\paragraph{A common-row fallback.}
For row functionals $F_i(f,g)$ before their fibre supremum,
$\sup_{f,g}\sum w_iF_i\le\sum w_i\sup_{f,g}F_i$.
Since the pre-supremum functional is affine in $(M,L,\delta a)$, its common
potential uses $\overline M=\sum w_iM_i$, $\overline L=\sum w_iL_i$ and
$\overline A=\sum w_i\delta_i a_i$; its dual subtracts
$m^{\mathsf T}\overline M^{-1}m$. Averaging $M_i^{-1}$ would be incorrect.
This replacement preserves the existing common-square mixing identity.
It cannot also claim the old full discarded-square deficit a second time.

\paragraph{Quadratic features close the missing first-chaos map exactly.}
The free completion in \eqref{eq:shared-stability-completion} is unnecessary
for an important feature class. For a Hermitian matrix $M$, let
\[
 A_M(z)=z^*Mz-\operatorname{tr}M,\qquad
 W_{jl;i}={:}z_jz_l\overline z_i{:}.
\]
The pointwise Wick identity is
\begin{equation}\label{eq:shared-quadratic-wick}
 A_Mz_l=\sum_jM_{lj}z_j+\sum_{i,j}M_{ij}W_{jl;i}.
\end{equation}
Consequently every disc-valued response, not only a near-maximizer, satisfies
\begin{align}
 [P(fA_M)]_l
 &=\sum_j\overline{M_{lj}}\langle z_j,f\rangle
       +\sum_{i,j}\overline{M_{ij}}\langle W_{jl;i},f\rangle
       \nonumber\\
 &=\sum_j\overline{M_{lj}}\langle z_j,Tf\rangle
       -\varepsilon^{-1}\sum_{i,j}\overline{M_{ij}}
                                      \langle W_{jl;i},Tf\rangle.
       \label{eq:shared-quadratic-closure}
\end{align}
Here $[Ph]_l=\langle z_l,h\rangle$. Reality of $A_M$ moves it to the test
probe, and \eqref{eq:shared-quadratic-wick} lies entirely in the first and
$(2,1)$ chaoses. Thus higher response coefficients vanish from the pairing.
No multiplication operator on all of $L^2$ is assumed: the bounded response
times $A_M$ is in $L^2$, and the probe product is a fixed polynomial.

Write the last expression as $\mathcal D_M(Tf)_l$. The full cross block is
now exactly
\[
 C_{ij}=\operatorname{Re}\langle\mathcal D_{M_i}(Tf),
                                      \mathcal D_{N_j}(Tg)\rangle.
\]
It is a contraction of the actual linear and cubic tensors. Changing its
unmeasured part independently while leaving those images fixed would violate
\eqref{eq:shared-quadratic-closure}. This addresses precisely the completion
freedom exposed by the score-feature diagnostic; it does not assert that
every quadratic-feature stability matrix is violated at a binding state.

The features are unbounded, so the bounded-feature remainder constant cannot
be reused. There is an explicit dimension-free substitute. If
$\|A_M\|_2=\|M\|_{\rm HS}=a$, then
\[
 \mathbb EA_M^4\le9a^4,\qquad
 \mathbb EA_M^6\le265a^6,\qquad
 \mathbb EA_M^8\le14833a^8.
\]
To prove these, diagonalize $M$ and write
$A_M=\sum\lambda_j(|z_j|^2-1)$. For $k\ge2$ its cumulant is
$(k-1)!\sum\lambda_j^k$, whose absolute value is at most $(k-1)!a^k$.
The moment--cumulant partition formula has nonnegative coefficients; replacing
each cumulant by this bound gives the corresponding even moment of
$a(\operatorname{Exp}(1)-1)$. Those centered even moments are the derangement
numbers $9,265,14833$.

For two such features with norms $a,b$, the $L^8$ remainder in
\S\ref{sec:fresh-second-chaos-stability}, using $\|Tf\|_2,\|Tg\|_2\le1$,
is at most
\[
 C(A,B)\le\frac{183}{4}(a^4+b^4)+\frac{45}{2}a^2b^2
                   +\frac{489}{20}(a^3b+ab^3)
             \le48(a^2+b^2)^2.
\]
Here $\sqrt{14833}<122$, $\sqrt{265}<163/10$ and $\varepsilon\le1$.
For the last inequality put $S=a^2+b^2$, $r=ab$; the difference multiplied
by $276$ is
$(138r-489S/20)^2+(9279/400)S^2\ge0$.
Therefore orthonormal Hermitian feature families give the necessary matrix
inequality
\begin{equation}\label{eq:shared-quadratic-stability}
 \boxed{\begin{pmatrix}D_f&\varepsilon C\\
                         \varepsilon C^{\mathsf T}&D_g\end{pmatrix}
             \succeq-16\sqrt{3\delta}\,I}
\end{equation}
for global $\delta$-near-maximizers. This also holds with the original signed
weights. No clipping error or free residual Gram completion is present. The
pointwise Wick identity, Hilbert-space probe closure and rational remainder
budget are kernel-checked in \code{QuadraticFeatures.lean}; the Gaussian
moment--cumulant application and matrix assembly here are written proofs.

The constant feature may be included without losing closure. For
$A=c+A_M$ with $\|A\|_2=a$, Minkowski and Cauchy--Schwarz give
$\|A\|_4\le2a$, $\|A\|_6\le\sqrt8\,a$ and
$\|A\|_8\le\sqrt{13}\,a$. Indeed $9=(\sqrt3)^4$,
$265<343=(\sqrt7)^6$ and $14833<20736=(\sqrt{12})^8$.
Since $(\sqrt8)^3<23$, the remainder becomes at most
\[
 \frac{507}{8}(a^4+b^4)+40a^2b^2+
                         \frac{69}{2}(a^3b+ab^3)
                 \le67(a^2+b^2)^2.
\]
For the last inequality, multiplying the difference by $347$ gives
$(347r/2-69S/2)^2+(541/8)S^2\ge0$.
Thus affine-quadratic feature families have matrix error at most
$4\sqrt{67\delta}\,I$. This second rational budget is also in Lean.

Thus the next uniform certificate can retain joint polynomial images and
their explicit quadratic contractions, use the exact local two-disc SDP for
moment prices, and bound the weighted image integrals. It still needs a
global joint-tensor inequality and complete row cover. A marginal radial law
alone contains neither the contraction map nor those weighted integrals.

\paragraph{Formalization and diagnostic status.}
The Hilbert-space Gram interval, its nested-projection specialization,
reflection estimate, scalar completion screen, factored prices, finite
moment-support assembly, pointwise alignment and cutoff inequalities,
cross-inner-product perturbation, scalar disc support and its attainment,
and near-maximizer price accounting are formalized in
\code{Grothe/LowerAnalytic/ResponseGram.lean}. The Gaussian application,
Gaussian image-lift and matrix-norm error estimates, the two-disc support
construction, matrix dilation and bilinear operator
identity above are paper proofs; they are not asserted to have been formalized
by those abstract lemmas. Discovery scripts and independent-grid results are
recorded in \code{docs/JOINT\_RESPONSE\_LIFT\_2026-10-23.md}.
No numerical endpoint is changed.

\code{TwoDisc.lean} additionally proves the positive-slack factorization and
its universal support bound, and sharp bounds and attainment for the four local
examples. These proofs use rational sum-of-squares factorizations, so the four
sharp local support values, unlike the floating SDP proposals, require no
external numerical hypothesis. The general rank-one SDP exactness argument
and the principal-minor proposal checker remain paper and external checks.
